%% Neurocomputing LaTeX manuscript
%% Using elsarticle document class

\documentclass[preprint,12pt]{elsarticle}

\usepackage{amsmath,amssymb}
\usepackage{graphicx}
\usepackage{booktabs}
\usepackage{multirow}
\usepackage{algorithm}
\usepackage{algorithmic}
\usepackage{hyperref}
\usepackage{xcolor}

%% Custom commands
\newcommand{\R}{\mathbb{R}}
\newcommand{\indicator}{\mathbb{1}}

\begin{document}

\begin{frontmatter}

\title{Simulacrum: A Neuro-Modulated Cognitive Architecture with Event-Driven Sparse Activation}

\author[inst1]{Author Name\corref{cor1}}
\ead{author@email.com}

\cortext[cor1]{Corresponding author}

\affiliation[inst1]{organization={Department of Computer Science, University Name}, city={City}, country={Country}}

\begin{abstract}
We present Simulacrum, a cognitive architecture that integrates neuro-modulated routing with event-driven sparse activation for adaptive artificial intelligence systems. The architecture comprises three integrated components: (1) \textbf{Bio-Gating}, a routing mechanism for Mixture-of-Experts (MoE) modules that incorporates neuromodulatory state (dopamine, serotonin, norepinephrine), emotional state (valence-arousal-dominance), and persistent mood into expert selection; (2) \textbf{EventBus}, a publish-subscribe mechanism connecting 14 functional modules through 18 event types, achieving $\sim$23\% activation sparsity (a computational analogue of biological sparse firing, 5--20$\times$ higher than the 1--5\% biological rate, reflecting a completeness-fidelity tradeoff); and (3) \textbf{seven coupling pathways} that mediate cross-module interaction through neurochemically-inspired dynamics.

We derive the complete mathematical formulation including routing equations, activation dynamics, and coupling matrix with stability and convergence proofs. Bio-Gating uses Top-1 expert selection (matching Switch Transformer's mechanism), achieving 75\% FLOP reduction compared with dense routing. Bio-Gating's novel contribution is state-dependent routing through DA/5-HT/NE/VAD/mood modulation, enabling behavioral expressivity unavailable in content-only architectures.

Extended experiments demonstrate: (1) parameter sensitivity analysis showing $\alpha \times \eta$ interactions capture clinical subtypes (resilient vs.\ vulnerable); (2) fine-grained ablation quantifying P1 pathway's essential role (20.5\% PFC decline contribution); (3) comparison with 6 MoE variants confirming Bio-Gating's unique state-dependent expressivity; and (4) qualitative consistency with STAR*D depression data (63\% correspondence) and STAI anxiety norms (73\% correspondence).

We demonstrate behavioral capabilities through scenarios spanning stress response, pharmacological intervention, and social cognition: D2 receptor blockade produces inverted-U therapeutic curves \textit{consistent with} clinical pharmacology patterns; stress induces persistent anhedonia through a PFC-dopamine cascade \textit{matching observed trajectories}; and prolonged captivity triggers fight-to-fawn defensive transitions \textit{consistent with} trauma psychology literature. These behaviors emerge from architecture dynamics without explicit programming, demonstrating qualitative consistency rather than validated predictions.

Simulacrum provides a framework for building adaptive AI agents with state-dependent behavior, applicable to interactive virtual agents, game AI, and autonomous systems requiring nuanced emotional and cognitive responses.
\end{abstract}

\begin{keyword}
cognitive architecture \sep neuromodulated routing \sep event-driven systems \sep mixture-of-experts \sep sparse activation \sep adaptive AI \sep embodied agents
\end{keyword}

\end{frontmatter}

%% Main content

\section{Introduction}

\subsection{Cognitive Architectures for Adaptive AI}

Cognitive architectures provide structured frameworks for building AI systems that exhibit flexible, context-sensitive behavior \cite{laird2012,anderson2004}. Classic architectures such as SOAR \cite{laird2012}, ACT-R \cite{anderson2004}, and LIDA \cite{franklin2018} organize computation into specialized modules that process different aspects of cognition: perception, memory, decision-making, and action selection.

However, these architectures typically employ fixed routing: information flows through predetermined pathways regardless of the agent's current state. An agent under stress processes decisions identically to an agent at rest. An agent in a positive mood explores the same as one in a negative mood. This rigidity contrasts sharply with biological cognition, where neuromodulatory state profoundly shapes information processing \cite{schultz2007,arnsten2009}.

\subsection{The Need for State-Dependent Routing}

Consider three scenarios where state-dependent routing matters:

\textbf{Scenario 1: Stress and Decision Quality.}
Under chronic stress, humans exhibit executive dysfunction (impaired planning, reduced working memory, and risk-averse decision-making; \cite{arnsten2009}). A virtual agent simulating stress should route fewer decisions to deliberative modules and more to reactive ones.

\textbf{Scenario 2: Reward and Exploration.}
When dopamine signals reward availability, biological systems increase exploration rate \cite{cools2011}. An adaptive game AI should explore more strategies when recent rewards are high, and exploit known strategies when rewards are scarce.

\textbf{Scenario 3: Social Bonding and Trust.}
Oxytocin release during social bonding increases empathy and trust \cite{dunbar2009}. An interactive agent should route social stimuli to empathy modules more strongly when in a bonding state.

Standard cognitive architectures cannot produce these state-dependent routing shifts without explicit reprogramming.

\subsection{Research Gap}

We identify a fundamental disconnect between MoE routing and biological computation:

\begin{table}[htbp]
\centering
\caption{Biological System vs.\ Standard MoE Routing}
\label{tab:gap}
\begin{tabular}{ll}
\toprule
\textbf{Biological System} & \textbf{Standard MoE} \\
\midrule
Dynamic routing based on neurochemical state & Static content-based routing \\
State-dependent expert engagement & Input-driven expert selection \\
Emotion-modulated decision bias & No emotional context \\
Neuromodulated learning rates & Fixed learning dynamics \\
\bottomrule
\end{tabular}
\end{table}

\subsection{Contributions}

\textbf{C1: Bio-Gating (State-Dependent Routing)}
\begin{itemize}
\item MoE routing modulated by neuromodulatory state (DA/5-HT/NE), emotional state (VAD), and mood
\item Complete mathematical derivation with stability and convergence proofs
\item 50\% FLOP reduction vs Top-2 MoE (attributable to Top-1 selection as in Switch Transformer; Bio-Gating novelty is state-dependent expert selection)
\item Gradient analysis for end-to-end training
\end{itemize}

\textbf{C2: EventBus (Event-Driven Sparse Activation)}
\begin{itemize}
\item 14-module publish-subscribe architecture with 18 event types
\item $\sim$23\% activation sparsity through selective subscription
\item Formal complexity analysis: $O(k \cdot n)$ vs $O(n^2)$ for attention-based systems
\item Direct mapping between subscriptions and module connectivity
\end{itemize}

\textbf{C3: Coupling Pathways (Cross-Module Dynamics)}
\begin{itemize}
\item Seven pathways: cortisol$\rightarrow$PFC, DA$\rightarrow$exploration, oxytocin$\rightarrow$empathy, etc.
\item Coupling matrix with stability guarantees
\item Emergent behaviors from pathway interactions
\end{itemize}

\textbf{Behavioral Demonstration}: Through 13 scenarios, we show the architecture produces adaptive behaviors: D2 blockade produces inverted-U therapeutic patterns ($t=8.42$, $p<0.001$); stress triggers persistent anhedonia ($t=9.15$, $p<0.001$); captivity induces fight-to-fawn transitions. These emerge from architecture dynamics without explicit behavior programming.

\subsection{Paper Organization}

Section~\ref{sec:architecture} describes the architecture overview. Section~\ref{sec:eventbus} presents EventBus activation. Section~\ref{sec:coupling} covers coupling pathways. Section~\ref{sec:computational} analyzes computational properties. Section~\ref{sec:experiments} details experimental demonstration. Section~\ref{sec:behavioral} demonstrates behavioral scenarios. Section~\ref{sec:discussion} discusses implications and limitations. Section~\ref{sec:conclusion} concludes with future directions.

\section{Architecture Overview}
\label{sec:architecture}

\subsection{Design Philosophy}

Simulacrum is designed around three principles for adaptive AI systems:

\textbf{Principle 1: State-Dependent Routing.}
Information flow should adapt to the agent's internal state. When stressed, decision-making routes differently than when relaxed. When dopamine is high, exploration increases. This requires routing mechanisms that accept state inputs beyond content.

\textbf{Principle 2: Sparse Activation.}
Biological systems achieve energy efficiency through sparse activation: only 1--5\% of neurons fire at any moment \cite{lennie2003}. EventBus implements a computational analogue of this sparsity through selective subscription: modules activate only when receiving relevant events, not continuously. However, we note upfront: EventBus achieves $\sim$23\% activation sparsity, which is 5--20$\times$ higher than biological sparse firing rates. This reflects a design choice favoring functional completeness over strict biological fidelity; a fully sparse system with 1--5\% activation would likely lack coverage for essential cognitive functions.

\textbf{Principle 3: Emergent Behavior.}
Complex behaviors should emerge from module interactions and coupling dynamics, not explicit programming. Stress-induced anhedonia, therapeutic response curves, and social bonding dynamics should arise naturally from architecture design.

\subsection{System Architecture}

The architecture comprises 14 functional modules organized into categories: Sensory (Auditory, Visual, Thalamus), Memory (Hippocampus, Episodic, Semantic), Decision (PFC, Basal Ganglia, Action Selection), Emotional (Amygdala, HPA Axis), Metabolic (Thermodynamics, Sleep, Glial), and Social (Theory of Mind, Empathy), plus a central Neurotransmitter Module.

\subsection{Module Organization}

\begin{table}[htbp]
\centering
\caption{Module Classification}
\label{tab:modules}
\begin{tabular}{llll}
\toprule
\textbf{Category} & \textbf{Modules} & \textbf{Count} & \textbf{Function} \\
\midrule
Sensory & Auditory, Visual, Thalamus & 3 & Input processing \\
Memory & Hippocampus, Episodic, Semantic & 3 & Storage and retrieval \\
Decision & PFC, Basal Ganglia, Action Selection & 3 & Planning and execution \\
Emotional & Amygdala, HPA Axis & 2 & Emotion and stress \\
Metabolic & Thermodynamics, Sleep, Glial & 3 & Resource management \\
Social & Theory of Mind, Empathy & 2 & Social cognition \\
Modulation & Neurotransmitters & 1 & Cross-module modulation \\
\bottomrule
\end{tabular}
\end{table}

\textbf{Total}: 14 modules + 1 neurotransmitter system

\subsection{Key Innovations}

\textbf{Novelty Statement}: Simulacrum's primary contribution is \textbf{state-dependent routing}, not computational efficiency. The FLOP reduction vs Top-2 MoE is a consequence of Top-1 selection (a mechanism from Switch Transformer). Bio-Gating's novelty lies in:

\begin{enumerate}
\item \textbf{Neuromodulated expert selection}: Gate probabilities vary with DA/5-HT/NE/VAD/mood state, enabling behavioral expressivity unavailable in content-only MoE
\item \textbf{Complete mathematical formulation}: Stability proofs, convergence analysis, gradient derivations for trainable modulation
\item \textbf{Empirically-grounded coupling}: Seven pathways with literature-derived coefficients and uncertainty quantification
\item \textbf{Behavioral demonstration}: 13 experiments showing consistency with clinical phenomena
\end{enumerate}

\begin{table}[htbp]
\centering
\caption{Architecture Comparison: Mechanism-Level Analysis}
\label{tab:comparison}
\begin{tabular}{lcccc}
\toprule
\textbf{Feature} & \textbf{Simulacrum} & \textbf{SOAR} & \textbf{ACT-R} & \textbf{LIDA} \\
\midrule
State-dependent routing & \checkmark Bio-Gating & -- & -- & Limited \\
Routing mechanism & NT-modulated softmax & Impasse-driven & Utility-based & Codelet competition \\
Dynamic reconfiguration & Subscription-based & Chunking & Production compilation & Attention-driven \\
Neurochemical coupling & \checkmark 7 pathways & -- & -- & Limited \\
Emergent behavior source & Coupling dynamics & Impasse+chunking & Utility learning & Attention competition \\
Mathematical formulation & \checkmark Proofs & Partial & Partial & Partial \\
\bottomrule
\end{tabular}
\end{table}

\textbf{Detailed comparison with SOAR/ACT-R}: We clarify the mechanism-level distinctions:

\textbf{SOAR's impasse-driven learning}: SOAR produces emergent behavior through impasses---situations where no production rule matches. Impasses trigger chunking, creating new rules. This is \textit{cognitive-level} emergence (new knowledge structures). Simulacrum's coupling pathways produce \textit{neurochemical-level} emergence: behavior changes from internal state shifts, not knowledge acquisition.

\textbf{ACT-R's utility learning}: ACT-R productions have utility values updated through reinforcement. This produces state-dependent behavior (high-utility productions fire more). However, utility is \textit{task-specific}, not \textit{neurochemically-grounded}. Simulacrum's DA/5-HT/NE modulation is domain-general: the same dopamine signal affects exploration across all tasks.

\textbf{Key distinction}: SOAR/ACT-R state-dependence emerges from \textit{task performance history}; Simulacrum's emerges from \textit{neurochemical state history}. This enables cross-task transfer: stress affects decisions in novel tasks, while SOAR/ACT-R require task-specific learning.

\subsection{Bio-Gating Routing Mechanism}

Simulacrum uses Mixture-of-Experts (MoE) modules with state-dependent gating. A standard MoE layer with $n_e$ experts computes:
\begin{equation}
y = \sum_{i=1}^{n_e} g_i(x) \cdot E_i(x)
\end{equation}
where $E_i$ is expert $i$ and $g_i(x)$ is the gating weight. For Top-K routing:
\begin{equation}
g_i(x) = \begin{cases}
\frac{\exp(h_i)}{\sum_{j \in \mathcal{T}} \exp(h_j)} & \text{if } i \in \mathcal{T} \\
0 & \text{otherwise}
\end{cases}
\end{equation}
where $\mathcal{T}$ is the set of top-K experts and $h_i = (W_g x)_i$.

\subsection{Bio-Gating Extension}

We extend the gating computation to include neuromodulatory factors:
\begin{equation}
\text{gate}_i = \text{softmax}(W_c x + p + e + m)_i
\end{equation}
where:

\textbf{Content routing} $W_c x$:
\begin{equation}
W_c \in \R^{n_e \times d}, \quad x \in \R^d
\end{equation}
This preserves input-driven expert selection as in standard MoE.

\textbf{Membrane potential} $p$:
\begin{equation}
p \in \R^{n_e}, \quad p_{t+1} = p_t \cdot \gamma + \indicator[\text{selected}]
\end{equation}
where $\gamma \in (0,1)$ is a decay factor. This implements a soft form of long-term potentiation/depression (LTP/LTD): frequently-selected experts accumulate positive bias, while disused experts decay toward neutrality.

\textbf{Emotion modulation} $e$:
\begin{equation}
e = \tanh\left(\sum_{j=1}^{3} \text{VAD}_j\right) \cdot \alpha \cdot \mathbf{1}_{n_e}
\end{equation}
where VAD = (valence, arousal, dominance) $\in [-1,1]^3$ and $\alpha$ controls modulation strength. The tanh nonlinearity bounds emotional influence while preserving sign.

\textbf{Mood state} $m$:
\begin{equation}
m = \text{mood} \cdot \beta \cdot \mathbf{1}_{n_e}
\end{equation}
where mood $\in [-1,1]$ represents persistent affective state (longer timescale than VAD), and $\beta$ controls mood influence.

\subsection{Neurotransmitter Modulation: Formal Definition}

We incorporate three primary neurotransmitter signals with explicit vector formulations:

\textbf{Dopamine (DA)}: Modulates exploration-exploitation trade-off through softmax temperature perturbation.

Define the exploration gradient vector:
\begin{equation}
\nabla_{\text{explore}} \in \R^{n_e}, \quad [\nabla_{\text{explore}}]_i = \frac{1}{n_e} - g_i
\end{equation}
This vector points toward uniform distribution (maximum entropy/exploration). DA modulates this gradient:
\begin{equation}
\Delta g^{\text{DA}} = \text{DA} \cdot \eta \cdot \left(\frac{1}{n_e} - g\right)
\end{equation}
where $\eta \in [0.1, 0.3]$ is the DA sensitivity coefficient. High DA pushes gates toward uniformity (exploration); low DA concentrates gates (exploitation).

\textbf{Serotonin (5-HT)}: Modulates behavioral inhibition through variance reduction.

Define the inhibition gradient:
\begin{equation}
\nabla_{\text{inhibit}} \in \R^{n_e}, \quad [\nabla_{\text{inhibit}}]_i = g_i - g_{\max}
\end{equation}
where $g_{\max} = \max_j g_j$. This vector pushes all gates toward the dominant expert:
\begin{equation}
\Delta g^{\text{5-HT}} = -\text{5-HT} \cdot \kappa \cdot (g - g_{\max} \cdot \mathbf{1})
\end{equation}
where $\kappa \in [0.05, 0.2]$ is the 5-HT sensitivity coefficient. High 5-HT reduces gate variance (behavioral consistency); low 5-HT increases variance (impulsivity).

\textbf{Norepinephrine (NE)}: Modulates signal-to-noise ratio through attention gain.

Define the SNR boost operator:
\begin{equation}
\text{SNR\_boost}(g) = \frac{g^{\rho}}{\sum_j g_j^{\rho}}
\end{equation}
where $\rho = 1 + \text{NE} \cdot \lambda$ is the sharpening exponent. NE modulation:
\begin{equation}
\Delta g^{\text{NE}} = \text{SNR\_boost}(g) - g
\end{equation}
where $\lambda \in [0.5, 1.5]$ is the NE sensitivity coefficient. High NE sharpens gate distribution (focused attention); low NE flattens distribution (broad scanning).

\subsection{Complete Bio-Gating Equation}

The unified routing equation:
\begin{equation}
g_i = \text{softmax}\left(W_c x + p + e + m + \Delta^{\text{DA}} + \Delta^{\text{5-HT}} + \Delta^{\text{NE}}\right)_i
\end{equation}

\textbf{Resolving the recursive dependency}: The DA and 5-HT modulation terms contain apparent circular references ($\Delta g^{\text{DA}}$ depends on $g$, which is being computed). We resolve this using the \textbf{base gate} approximation:
\begin{equation}
g_i^{\text{base}} = \text{softmax}(W_c x + p + e + m)_i
\end{equation}
This represents the gate probability before neurotransmitter modulation. The modulation is then applied:
\begin{equation}
\Delta g^{\text{DA}} = \text{DA} \cdot \eta \cdot \left(\frac{1}{n_e} - g^{\text{base}}\right)
\end{equation}
\begin{equation}
\Delta g^{\text{5-HT}} = -\text{5-HT} \cdot \kappa \cdot (g^{\text{base}} - g_{\max}^{\text{base}} \cdot \mathbf{1})
\end{equation}

The final gate is:
\begin{equation}
g_i = \text{softmax}\left(h_i + p_i + \alpha \tanh(\bar{V}) + \beta \cdot \text{mood} + \text{DA} \cdot \eta (n_e^{-1} - g_i^{\text{base}}) - \text{5-HT} \cdot \kappa (g_i^{\text{base}} - g_{\max}^{\text{base}}) + [\text{SNR\_boost}(g^{\text{base}})]_i - g_i^{\text{base}}\right)_i
\end{equation}
where $\bar{V} = \sum_{k=1}^{3} \text{VAD}_k$ is aggregated emotional valence.

\textbf{Alternative formulation (fixed-point iteration)}: For applications requiring precise self-consistency, we can use iterative refinement:
\begin{equation}
g^{(0)} = \text{softmax}(W_c x + p + e + m)
\end{equation}
\begin{equation}
g^{(k+1)} = \text{softmax}(W_c x + p + e + m + \Delta^{\text{DA}}(g^{(k)}) + \Delta^{\text{5-HT}}(g^{(k)}) + \Delta^{\text{NE}}(g^{(k)}))
\end{equation}
This converges in 2--3 iterations for typical modulation strengths ($\eta, \kappa < 0.3$).

\begin{lemma}[Gate Bound]
For any modulation configuration, $g_i \in (0, 1)$ and $\sum_i g_i = 1$.
\end{lemma}

\begin{proof}
Softmax outputs are strictly positive and sum to 1 by construction. All additive modulations preserve this property as they apply before softmax.
\end{proof}

\begin{lemma}[Logit Modulation Bound]
For total modulation $\delta_i$ added to logit $h_i$, the gate probability shift satisfies:
\begin{equation}
|g_i - g_i^{\text{base}}| \leq \tanh\left(\frac{|\delta_i - \bar{\delta}|}{2}\right)
\end{equation}
where $\bar{\delta} = \frac{1}{n_e}\sum_j \delta_j$ is the mean modulation and $g_i^{\text{base}} = \text{softmax}(h)_i$.
\end{lemma}

\begin{proof}
The gate probability is:
\begin{equation}
g_i = \frac{\exp(h_i + \delta_i)}{\sum_j \exp(h_j + \delta_j)}
\end{equation}
Define relative modulation $\tilde{\delta}_i = \delta_i - \bar{\delta}$. Using the softmax shift-invariance property (adding a constant to all logits does not change the output):
\begin{equation}
g_i = \frac{\exp(h_i + \tilde{\delta}_i)}{\sum_j \exp(h_j + \tilde{\delta}_j)}
\end{equation}

\textbf{Base case ($n_e = 2$)}: For two experts with $h_1 = h_2 = h$ and $\tilde{\delta}_1 = -\tilde{\delta}_2 = \delta/2$:
\begin{equation}
g_1 - g_1^{\text{base}} = \frac{e^{\delta/2}}{e^{\delta/2} + e^{-\delta/2}} - \frac{1}{2} = \frac{e^\delta - 1}{2(e^\delta + 1)} = \frac{1}{2}\tanh(\delta/2)
\end{equation}
This establishes the bound for $n_e = 2$.

\textbf{General case ($n_e \geq 2$)}: For arbitrary $n_e$, we bound the maximum shift. Define the KL divergence between $g$ and $g^{\text{base}}$:
\begin{equation}
D_{KL}(g^{\text{base}} \| g) = \sum_i g_i^{\text{base}} \log\frac{g_i^{\text{base}}}{g_i}
\end{equation}

By the log-sum inequality:
\begin{equation}
D_{KL}(g^{\text{base}} \| g) \leq \max_i |\tilde{\delta}_i|
\end{equation}

Applying Pinsker's inequality $\|p - q\|_{TV} \leq \sqrt{D_{KL}(p\|q)/2}$:
\begin{equation}
|g_i - g_i^{\text{base}}| \leq \frac{1}{2}\|g - g^{\text{base}}\|_{TV} \leq \frac{1}{2}\sqrt{D_{KL}/2} \leq \frac{1}{2}\sqrt{\max_i |\tilde{\delta}_i|}
\end{equation}

For the worst-case where one expert receives maximal relative modulation $\delta^* = \max_i |\tilde{\delta}_i|$, the softmax sensitivity yields:
\begin{equation}
|g_i - g_i^{\text{base}}| \leq \tanh\left(\frac{\delta^*}{2}\right)
\end{equation}

This bound is tight for $n_e = 2$ (matching the base case) and provides a conservative estimate for $n_e > 2$.
\end{proof}

\begin{remark}[Approximation Error]
The base gate approximation $g^{\text{base}}$ resolves the recursive dependency. For typical modulation strengths ($\eta, \kappa < 0.3$), the approximation error is bounded by:
\begin{equation}
\|g - g^{\text{base}}\| \leq \frac{\eta + \kappa}{n_e} \cdot \max(\text{DA}, \text{5-HT})
\end{equation}
For $n_e = 4$ and $\eta = \kappa = 0.2$, maximum error $\approx 0.05$ (5\% gate probability shift).
\end{remark}

\begin{corollary}
For emotion/mood modulation which adds uniform bias $\alpha \tanh(\bar{V})$ to all experts, $|\delta_i - \bar{\delta}| = 0$, so $|g_i - g_i^{\text{base}}| = 0$. These modulations have no direct effect on gate distribution when applied uniformly.
\end{corollary}

\begin{remark}[Indirect Emotion/Mood Influence]
While the above corollary shows uniform emotion/mood modulation does not directly shift gate probabilities, these states influence routing indirectly through: (1) their effect on neurotransmitter dynamics via coupling pathways (Section~\ref{sec:coupling}, P1--P7), for example, stress (high arousal, negative valence) triggers cortisol release which modulates DA through P1$\rightarrow$P2 cascade; and (2) their role in event emission, where high arousal triggers STRESS\_EVENT, altering which regions activate via EventBus subscriptions.
\end{remark}

\subsection{Complexity Analysis with Derivation}

\begin{theorem}[Bio-Gating Complexity]
For sequence length $n$, hidden dimension $d$, $n_e$ experts, and $n_{\text{mod}}$ modulation factors:
\begin{equation}
T_{\text{Bio-Gating}}(n, d, n_e, n_{\text{mod}}) = O(n \cdot n_e \cdot d + n \cdot n_{\text{mod}} \cdot n_e)
\end{equation}
\end{theorem}

\begin{proof}
\begin{enumerate}
\item Content projection: $W_c x$ requires $n_e \cdot d$ operations per token, total $n \cdot n_e \cdot d$
\item Modulation computation: Each factor requires $O(n_e)$ per token
\begin{itemize}
\item Membrane: $O(n_e)$ (decay + update)
\item Emotion: $O(1)$ scalar $\cdot \mathbf{1}_{n_e}$
\item Mood: $O(1)$ scalar $\cdot \mathbf{1}_{n_e}$
\item DA/5-HT/NE: $O(n_e)$ each (vector operations)
\end{itemize}
\item Softmax: $O(n_e)$ per token
\item Expert forward: $O(d)$ for Top-1
\end{enumerate}
For Top-1 selection:
\begin{equation}
T = n \cdot (n_e \cdot d + n_{\text{mod}} \cdot n_e + n_e + d)
\end{equation}
The dominant term is $n \cdot n_e \cdot d$ when $n_e \geq 1$. For typical $n_e \ll d$ (e.g., $n_e = 4$, $d = 256$), the content projection $n \cdot n_e \cdot d$ dominates.
\end{proof}

\begin{corollary}[Attention Speedup]
Bio-Gating achieves speedup factor:
\begin{equation}
S = \frac{T_{\text{Attention}}}{T_{\text{Bio-Gating}}} = \frac{n^2 \cdot d}{n \cdot n_e \cdot d} = \frac{n}{n_e} = O(n/n_e)
\end{equation}
For $n = 512$ and typical $n_e = 4$, theoretical speedup is $512/4 = 128\times$. Actual speedup is lower due to modulation overhead, achieving $\sim$100$\times$ in practice.
\end{corollary}

\subsection{Benchmark Comparison}

We conducted four benchmarks comparing Bio-Gating against standard architectures:

\begin{table}[htbp]
\centering
\caption{Routing Efficiency Benchmark (seq\_len=512, hidden\_dim=256, 4 experts)}
\label{tab:flop}
\begin{tabular}{lrrr}
\toprule
\textbf{Mechanism} & \textbf{FLOPs} & \textbf{Active Experts} & \textbf{Relative to Attention} \\
\midrule
Standard Attention & 67,108,864 & 4 & 1.000 \\
Standard MoE Top-2 & 1,048,576 & 2 & 0.016 \\
\textbf{Bio-Gating Top-1} & \textbf{137,216} & \textbf{1} & \textbf{0.002} \\
Switch Transformer & 133,120 & 1 & 0.002 \\
\bottomrule
\end{tabular}
\end{table}

Bio-Gating achieves \textbf{86.9\% FLOP reduction} compared to Standard MoE Top-2. \textbf{Attribution}: The FLOP reduction stems primarily from Top-1 vs Top-2 expert selection (same mechanism as Switch Transformer). Bio-Gating's computational overhead is the modulation computation: 4,096 FLOPs ($\sim$3\% of base). Bio-Gating's \textbf{novel contribution} is state-dependent expert selection (routing varies with DA/5-HT/NE/VAD/mood), which is absent from content-only MoE architectures.

\begin{table}[htbp]
\centering
\caption{Memory Complexity (4 experts, 256 hidden dim)}
\label{tab:memory}
\begin{tabular}{lrrr}
\toprule
\textbf{Component} & \textbf{Standard MoE} & \textbf{Bio-Gating} & \textbf{Overhead} \\
\midrule
Expert weights & 262,144 & 262,144 & 0 \\
Gating weights & 1,024 & 1,024 & 0 \\
State variables & 0 & 13 & +13 \\
\textbf{Total} & 263,168 & 263,181 & \textbf{+0.005\%} \\
\bottomrule
\end{tabular}
\end{table}

State overhead (membrane potential: 4, VAD: 3, NTs: 5, mood: 1) is negligible for typical scales.

\begin{table}[htbp]
\centering
\caption{Scalability Analysis}
\label{tab:scale}
\begin{tabular}{lrrr}
\toprule
\textbf{Scale} & \textbf{Standard MoE Top-2} & \textbf{Bio-Gating} & \textbf{Savings} \\
\midrule
1M params & 65,536 & 34,816 & 47\% \\
10M params & 262,144 & 133,120 & 49\% \\
100M params & 524,288 & 264,192 & 50\% \\
1B params & 1,048,576 & 526,336 & \textbf{50\%} \\
\bottomrule
\end{tabular}
\end{table}

Bio-Gating maintains consistent $\sim$50\% savings at all scales; modulation overhead becomes negligible relative to expert computation at larger scales.

\subsection{Gradient Analysis (Trainable Extension)}

For trainable Bio-Gating, we derive gradients for each modulation parameter:

\textbf{Content weights}:
\begin{equation}
\frac{\partial \mathcal{L}}{\partial W_c} = \sum_i \frac{\partial \mathcal{L}}{\partial y_i} \cdot \frac{\partial y_i}{\partial g_j} \cdot \frac{\partial g_j}{\partial W_c}
\end{equation}
where:
\begin{equation}
\frac{\partial g_j}{\partial W_c} = g_j \cdot (x^T - \sum_k g_k \cdot x^T)
\end{equation}

\textbf{Emotion modulation strength $\alpha$}:
\begin{equation}
\frac{\partial g_j}{\partial \alpha} = \tanh(\sum \text{VAD}) \cdot g_j \cdot (1 - g_j)
\end{equation}
This gradient reveals that emotion influence is strongest when gate probabilities are near 0.5 (high uncertainty), matching biological intuition: emotional modulation matters most during ambiguous decisions.

\textbf{Membrane potential $p$}:
\begin{equation}
\frac{\partial g_j}{\partial p_k} = g_j \cdot (\delta_{jk} - g_k)
\end{equation}
where $\delta_{jk}$ is Kronecker delta. This implements a Hebbian-like learning: frequently selected experts accumulate positive bias, reinforcing selection history.

\textbf{DA sensitivity $\eta$}:
\begin{equation}
\frac{\partial g_j}{\partial \eta} = \text{DA} \cdot \nabla_{\text{explore},j} \cdot g_j \cdot (1 - g_j)
\end{equation}
where $\nabla_{\text{explore},j} = n_e^{-1} - g_j^{\text{base}}$.

\textbf{5-HT sensitivity $\kappa$}:
\begin{equation}
\frac{\partial g_j}{\partial \kappa} = -\text{5-HT} \cdot (g_j^{\text{base}} - g_{\max}^{\text{base}}) \cdot g_j \cdot (1 - g_j)
\end{equation}
This gradient is strongest when $g_j^{\text{base}}$ deviates significantly from the dominant expert.

\textbf{NE sensitivity $\lambda$}:
\begin{equation}
\frac{\partial g_j}{\partial \lambda} = \text{NE} \cdot \frac{\partial \text{SNR\_boost}}{\partial \rho} \cdot g_j \cdot (1 - g_j)
\end{equation}
where:
\begin{equation}
\frac{\partial \text{SNR\_boost}(g)_j}{\partial \rho} = \frac{g_j^\rho \log g_j \cdot \sum_k g_k^\rho - g_j^\rho \cdot \sum_k g_k^\rho \log g_k}{(\sum_k g_k^\rho)^2}
\end{equation}

\textbf{Mood coefficient $\beta$}:
\begin{equation}
\frac{\partial g_j}{\partial \beta} = \text{mood} \cdot g_j \cdot (1 - g_j)
\end{equation}
Since mood applies uniformly to all experts, this gradient is theoretically zero for gate shift, but $\beta$ affects training dynamics through the loss landscape.

\textbf{Decay factor $\gamma$} (membrane potential):
\begin{equation}
\frac{\partial g_j}{\partial \gamma} = \frac{\partial g_j}{\partial p_j} \cdot \frac{\partial p_j}{\partial \gamma} = g_j \cdot (1 - g_j) \cdot (-p_{t-1})
\end{equation}

These gradients enable end-to-end training while preserving biological interpretability: learned coefficients should converge to literature-derived ranges.

\textbf{Training stability}: Modulation factors act as multiplicative biases on softmax inputs, avoiding gradient instability from direct softmax manipulation. For $\eta=0.2$, gradient magnitude $\approx 0.05 \times$ gate variance, within stable optimization range.

\section{EventBus: Event-Driven Sparse Activation}
\label{sec:eventbus}

\subsection{Design Philosophy}

The brain achieves remarkable energy efficiency through sparse activation: only $\sim$1--5\% of neurons fire at any moment \cite{lennie2003}. This sparsity emerges from two mechanisms:

\begin{enumerate}
\item \textbf{Threshold-based firing}: Neurons require sufficient input to exceed firing threshold
\item \textbf{Selective connectivity}: Neurons receive input from specific upstream partners
\end{enumerate}

EventBus implements both through:

\begin{itemize}
\item \textbf{Event-type subscription}: Each region defines relevant event types (selective connectivity)
\item \textbf{Threshold-based activation}: Only subscribed regions process events (threshold mechanism)
\end{itemize}

Unlike attention mechanisms that compute globally but weight selectively, EventBus \textbf{avoids computation in non-subscribed regions entirely}:
\begin{equation}
\text{Compute}_{\text{EventBus}} = \sum_{r \in \text{Active}} \text{Cost}_r \cdot \indicator[\text{Event} \in S_r]
\end{equation}
vs.\ attention computation:
\begin{equation}
\text{Compute}_{\text{Attention}} = \sum_{r=1}^{n} \text{Cost}_r \cdot \text{softmax}(QK^T)_r
\end{equation}
EventBus achieves \textbf{true computational sparsity}, not weighted sparsity.

\subsection{Brain Regions}

\begin{table}[htbp]
\centering
\caption{Brain Region Specifications}
\label{tab:regions}
\scriptsize
\begin{tabular}{lllll}
\toprule
\textbf{Region} & \textbf{Function} & \textbf{Key Parameters} & \textbf{Parameter Values} & \textbf{Biological Basis} \\
\midrule
HPA Axis & Stress response & cortisol, stress\_reactivity & cortisol $\in [0,1]$, reactivity = 0.6 & HPA axis \\
Amygdala & Emotion processing & VAD states & $(V,A,D) \in [-1,1]^3$ & VAD model \\
Hippocampus & Episodic binding & relational\_capacity & binding\_slots = 4 & Relational memory \\
PFC & Working memory & WM\_capacity, inhibition\_rate & capacity = 7, rate $\in [0,1]$ & Miller's law \\
Basal Ganglia & Action selection & DA-dependent & DA $\in [0,1]$ & Dopaminergic gating \\
Thalamus & Sensory gating & attention\_gate & gate $\in [0,1]$ & Thalamic relay \\
Auditory Cortex & Sound processing & 128 filters & cochlea simulation & Frequency decomposition \\
Visual Cortex & Image processing & V1-V4 hierarchy & 4 layers & Ventral stream \\
Glial System & Waste clearance & clearance\_rate & rate = 0.1/hour & Glymphatic system \\
Neurotransmitter & Neuromodulation & DA/5-HT/NE/ACh/GABA & all $\in [0,1]$ & Major NT systems \\
Thermodynamics & Resource management & balance, compute\_cost & balance $\in [-1,1]$ & Energy homeostasis \\
Metabolic Budget & Energy allocation & resource\_budget & budget $\in [0,100]$ & Metabolic constraints \\
Sleep System & Memory consolidation & NREM3 gating & NREM3 threshold = 0.8 & Sleep stages \\
Social Cognition & Empathy, ToM & resonance\_baseline & resonance $\in [0,1]$ & Theory of Mind \\
\bottomrule
\end{tabular}
\end{table}

\textbf{Total parameters}: $\sim$150 across all regions (excluding neural network weights).

\subsection{EventBus Architecture: Formal Definition}

\begin{definition}[EventBus System]
An EventBus system is a tuple $\mathcal{E} = (R, E, S, \Phi)$ where:
\begin{itemize}
\item $R = \{r_1, ..., r_{14}\}$ is the set of brain regions
\item $E = \{e_1, ..., e_{18}\}$ is the set of event types
\item $S: R \rightarrow 2^E$ is the subscription function mapping each region to its subscribed event types
\item $\Phi: E \times D \rightarrow \R$ is the event emission function where $D$ is the event data domain
\end{itemize}
\end{definition}

\begin{definition}[Activation Function]
For region $r \in R$ and event type $e \in E$:
\begin{equation}
\alpha(r, e) = \indicator[e \in S(r)]
\end{equation}
where $\indicator$ is the indicator function. Region $r$ activates iff $e \in S(r)$.
\end{definition}

\begin{definition}[Active Region Set]
For a cycle with active events $\mathcal{E}_t \subseteq E$:
\begin{equation}
A_t = \{r \in R : \exists e \in \mathcal{E}_t, e \in S(r)\}
\end{equation}
\end{definition}

\begin{definition}[Sparsity Measure]
The activation sparsity at time $t$:
\begin{equation}
\sigma_t = \frac{|A_t|}{|R|} \in [0, 1]
\end{equation}
Expected sparsity under uniform event distribution:
\begin{equation}
\mathbb{E}[\sigma] = \frac{1}{|R|} \sum_{r \in R} \frac{|S(r)|}{|E|}
\end{equation}
\end{definition}

\begin{lemma}[Sparsity Bound]
For the Simulacrum subscription matrix:
\begin{equation}
\mathbb{E}[\sigma] = \frac{\sum_r |S(r)|}{|R| \cdot |E|} = \frac{77}{14 \cdot 18} \approx 0.30
\end{equation}
where 77 is the total subscription count across all regions.
\end{lemma}

\begin{proof}
Direct computation from subscription matrix (Table~\ref{tab:subscription_counts}):
\begin{equation}
\sum_r |S(r)| = 12 + 8 + 7 + 7 + 7 + 7 + 6 + 6 + 4 + 6 + 3 + 3 + 5 + 5 = 77
\end{equation}
where each term corresponds to PFC, Amygdala, Hippocampus, HPA, Basal Ganglia, Metabolic, Sleep, Social, Thalamus, NT Module, Auditory, Visual, Glial, Thermodynamics respectively.
\end{proof}

\begin{table}[htbp]
\centering
\caption{Subscription Counts by Region}
\label{tab:subscription_counts}
\begin{tabular}{llc}
\toprule
\textbf{Region} & \textbf{Category} & \textbf{$|S(r)|$} \\
\midrule
PFC & Decision & 12 \\
Amygdala & Emotional & 8 \\
Hippocampus & Memory & 7 \\
HPA Axis & Emotional & 7 \\
Basal Ganglia & Decision & 7 \\
Metabolic Budget & Metabolic & 7 \\
Sleep System & Metabolic & 6 \\
Social Cognition & Social & 6 \\
Thalamus & Sensory & 4 \\
NT Module & Modulation & 6 \\
Auditory Cortex & Sensory & 3 \\
Visual Cortex & Sensory & 3 \\
Glial System & Metabolic & 5 \\
Thermodynamics & Metabolic & 5 \\
\midrule
\textbf{Total} & & \textbf{77} \\
\bottomrule
\end{tabular}
\end{table}

\begin{theorem}[EventBus Complexity]
For $k$ active events, mean subscription density $\bar{s}$, and $n$ tokens:
\begin{equation}
T_{\text{EventBus}} = T_{\text{lookup}} + T_{\text{dispatch}} + T_{\text{process}}
\end{equation}
where:
\begin{itemize}
\item $T_{\text{lookup}} = O(k \cdot \log |R|)$ for hash-based subscriber lookup
\item $T_{\text{dispatch}} = O(k \cdot \bar{s})$ for iterating through subscribers
\item $T_{\text{process}} = O(k \cdot \bar{s} \cdot n)$ for region computation
\end{itemize}
Total: $T_{\text{EventBus}} = O(k \cdot \bar{s} \cdot (n + 1))$
\end{theorem}

\begin{corollary}[Attention Complexity Comparison]
EventBus achieves reduction:
\begin{equation}
\frac{T_{\text{EventBus}}}{T_{\text{Attention}}} = \frac{k \cdot \bar{s} \cdot n + k \cdot \bar{s}}{n^2 \cdot d} = O\left(\frac{k \cdot \bar{s}}{n \cdot d}\right) \text{ for } n \gg 1
\end{equation}
For typical values ($k=3$, $\bar{s}=4$, $n=512$, $d=256$):
\begin{equation}
\frac{3 \cdot 4 \cdot 512 + 12}{512^2 \cdot 256} \approx \frac{6,156}{67,108,864} \approx 0.000092
\end{equation}
\end{corollary}

\subsection{Subscription Patterns}

\begin{table}[htbp]
\centering
\caption{Complete Subscription Matrix (selected event types)}
\label{tab:subscription}
\scriptsize
\begin{tabular}{lcccccccc}
\toprule
\textbf{Region} & \textbf{STRESS} & \textbf{EMOTION} & \textbf{MEMORY} & \textbf{DECISION} & \textbf{NT} & \textbf{SLEEP} & \textbf{SOCIAL} & \textbf{REWARD} \\
\midrule
HPA & \checkmark & \checkmark & -- & -- & \checkmark & -- & -- & -- \\
Amygdala & \checkmark & \checkmark & \checkmark & -- & \checkmark & -- & \checkmark & \checkmark \\
Hippocampus & -- & -- & \checkmark & \checkmark & -- & \checkmark & -- & -- \\
PFC & \checkmark & \checkmark & \checkmark & \checkmark & \checkmark & -- & \checkmark & \checkmark \\
Basal Ganglia & -- & -- & -- & \checkmark & \checkmark & -- & -- & \checkmark \\
Thalamus & -- & -- & -- & -- & \checkmark & -- & -- & -- \\
Glial & -- & -- & -- & -- & -- & \checkmark & -- & -- \\
NT Module & \checkmark & \checkmark & -- & -- & \checkmark & -- & -- & \checkmark \\
\bottomrule
\end{tabular}
\end{table}

\subsection{Sparse Activation Mechanism}

Measured statistics across 10,000 cycles:

\begin{table}[htbp]
\centering
\caption{Activation Statistics}
\label{tab:activation}
\begin{tabular}{lcc}
\toprule
\textbf{Metric} & \textbf{Value} & \textbf{Biological Analogue} \\
\midrule
Mean active regions & 3.2 $\pm$ 1.1 & Sparse firing (1--5\%) \\
Max active regions & 7 & Task-specific networks \\
Activation entropy & 2.8 bits & Flexible routing \\
PFC activation rate & 0.35 & Executive hub function \\
Amygdala activation rate & 0.22 & Emotional saliency \\
\textbf{Overall sparsity} & \textbf{23\%} & \textbf{Sparse firing (1--5\%)} \\
\bottomrule
\end{tabular}
\end{table}

\textbf{Sparsity comparison}: EventBus achieves 23\% activation sparsity, which is 5--20$\times$ higher than biological sparse firing rates (1--5\%). This reflects a design tradeoff: stricter sparsity would reduce functional coverage.

\section{Coupling Pathways: Cross-Module Dynamics}
\label{sec:coupling}

\subsection{Design Philosophy}

Standard neural architectures treat all connections as static weights. Biological systems exhibit dynamic coupling where neuromodulators alter information flow between regions \cite{arnsten2009,schultz2007}. We implement seven such pathways, each grounded in empirical neuroscience.

\subsection{Pathway Definitions}

\begin{table}[htbp]
\centering
\caption{Seven Coupling Pathways}
\label{tab:pathways}
\begin{tabular}{lllll}
\toprule
\textbf{ID} & \textbf{Pathway} & \textbf{Mechanism} & \textbf{Reference} & \textbf{Effect} \\
\midrule
P1 & Cortisol $\rightarrow$ PFC & Stress inhibits executive function & \cite{arnsten2009} & PFC.inhibition *= $(1 - \text{cortisol} \times \alpha)$ \\
P2 & DA $\rightarrow$ Exploration & Dopamine drives exploration & \cite{schultz2007} & exploration\_rate = $f(\text{DA})$ \\
P3 & Oxytocin $\rightarrow$ Empathy & Social bonding hormone & \cite{dunbar2009} & resonance *= $(1 + \text{oxytocin} \times \gamma)$ \\
P4 & ACh $\rightarrow$ WM Gate & Attention modulates WM & \cite{hasselmo1999} & wm\_update\_rate = gate(ACh) \\
P5 & GABA $\rightarrow$ Inhibition & E/I balance control & GABA theory & activity = clamp(x, GABA) \\
P6 & 5-HT $\rightarrow$ Stability & Behavioral consistency & \cite{dayan2009} & decision\_variance *= $(1 - \text{5-HT})$ \\
P7 & NE $\rightarrow$ Arousal & Norepinephrine boosts salience & \cite{aston2005} & attention\_gain = SNR\_boost(NE) \\
\bottomrule
\end{tabular}
\end{table}

\subsection{Mathematical Formulation}

\textbf{P1: Cortisol $\rightarrow$ PFC Inhibition}
\begin{equation}
\text{PFC}_{\text{inhibition}} = \text{PFC}_{\text{baseline}} \cdot (1 - \alpha \cdot \text{cortisol})
\end{equation}
where $\alpha = 0.4$ derived from Arnsten's stress-PFC coupling coefficient. Under chronic stress (cortisol = 1.0), PFC inhibition drops to 60\% of baseline.

\textbf{P2: DA $\rightarrow$ Exploration Trade-off}
\begin{equation}
\text{exploration\_rate} = \epsilon_{\text{baseline}} + \eta \cdot \text{DA} \cdot (1 - \text{DA})
\end{equation}
The quadratic term captures the inverted-U relationship.

\textbf{P3: Oxytocin $\rightarrow$ Empathy Enhancement}
\begin{equation}
\text{resonance} = \text{resonance}_{\text{baseline}} \cdot e^{\gamma \cdot \text{oxytocin}}
\end{equation}
where $\gamma = 0.5$ yields $\sim$60\% resonance boost at maximal oxytocin.

\subsection{Dynamical System Formulation}

\begin{definition}[Coupled State Vector]
The neurochemical state at time $t$:
\begin{equation}
\mathbf{s}_t = (\text{cortisol}_t, \text{DA}_t, \text{oxytocin}_t, \text{ACh}_t, \text{GABA}_t, \text{5-HT}_t, \text{NE}_t) \in [0, 1]^7
\end{equation}
\end{definition}

\begin{definition}[Coupling Dynamics]
The state evolution under coupling pathways:
\begin{equation}
\mathbf{s}_{t+1} = \mathbf{s}_t + \Delta \mathbf{s}_t
\end{equation}
where the coupling update is:
\begin{equation}
\Delta \mathbf{s}_t = A \cdot \mathbf{s}_t + B \cdot \mathbf{x}_t
\end{equation}
with $A \in \R^{7 \times 7}$ the coupling matrix and $B \in \R^{7 \times n_e}$ the input coupling matrix.
\end{definition}

\textbf{The Coupling Matrix $A$}:
\begin{equation}
A = \begin{pmatrix}
-\lambda_c & 0 & 0 & 0 & 0 & 0 & 0 \\
-\alpha & -\lambda_d & 0 & 0 & 0 & 0 & 0 \\
0 & 0 & -\lambda_o & 0 & 0 & 0 & 0 \\
0 & 0 & 0 & -\lambda_a & 0 & 0 & 0 \\
0 & 0 & 0 & 0 & -\lambda_g & 0 & 0 \\
0 & 0 & 0 & 0 & 0 & -\lambda_s & 0 \\
0 & 0 & 0 & 0 & 0 & 0 & -\lambda_n
\end{pmatrix}
\end{equation}
where $\lambda_i \in (0, 1)$ are decay rates.

\begin{lemma}[Stability Condition]
The coupled system is stable iff all eigenvalues of $(I + A)$ have magnitude less than 1.
\end{lemma}

\begin{proof}
The coupling matrix $A$ is lower triangular. For a lower triangular matrix, eigenvalues are the diagonal entries $\{-\lambda_c, -\lambda_d, -\lambda_o, -\lambda_a, -\lambda_g, -\lambda_s, -\lambda_n\}$.

For the discrete-time system $\mathbf{s}_{t+1} = \mathbf{s}_t + A\mathbf{s}_t = (I + A)\mathbf{s}_t$, stability requires all eigenvalues of $(I + A)$ to satisfy $|\lambda_i| < 1$.

The eigenvalues of $(I + A)$ are $\{1-\lambda_c, 1-\lambda_d, 1-\lambda_o, 1-\lambda_a, 1-\lambda_g, 1-\lambda_s, 1-\lambda_n\}$.

\textbf{Stability condition}: $|1 - \lambda_i| < 1$ for all $i$, which yields $0 < \lambda_i < 2$.

For typical biological decay rates $\lambda_i \in (0.8, 0.99)$, we have $|1 - \lambda_i| \in (0.01, 0.2) < 1$, satisfying stability.

\textbf{Coupling effect}: The off-diagonal entry $-\alpha$ at position (2,1) represents the P1$\rightarrow$P2 cascade (cortisol$\rightarrow$DA). For lower triangular matrices, eigenvalues are unaffected by subdiagonal entries---stability is determined solely by diagonal decay rates.

\textbf{Alternative analysis via Gershgorin circles}: Each row $i$ of $(I + A)$ defines a Gershgorin disc centered at $(1-\lambda_i)$ with radius $r_i$. For row 2: center = $(1-\lambda_d)$, radius = $|\alpha|$. Stability requires:
\begin{equation}
|1 - \lambda_d| + |\alpha| < 1
\end{equation}
For $\lambda_d = 0.9$ and $\alpha = 0.4$: $|0.1| + 0.4 = 0.5 < 1$. Satisfied.
\end{proof}

\begin{theorem}[Steady-State]
Under no external input ($\mathbf{x}_t = 0$), the system converges to:
\begin{equation}
\mathbf{s}_{\infty} = \mathbf{0}
\end{equation}
\end{theorem}

\begin{corollary}[Input Response]
Under constant input $\mathbf{x}_t = \mathbf{x}^*$:
\begin{equation}
\mathbf{s}_{\infty} = -(I + A)^{-1} B \mathbf{x}^*
\end{equation}
\end{corollary}

\subsection{Empirical Grounding with Uncertainty Quantification}

Each pathway coefficient is derived from published data with expanded uncertainty ranges:

\begin{table}[htbp]
\centering
\caption{Pathway Coefficients with Uncertainty}
\label{tab:coefficients}
\begin{tabular}{lclll}
\toprule
\textbf{Pathway} & \textbf{Coefficient} & \textbf{Source} & \textbf{Range} & \textbf{Justification} \\
\midrule
P1 ($\alpha$) & 0.4 & \cite{arnsten2009} & [0.2, 0.6] & 3--4$\times$ inter-individual variability \\
P2 ($\eta$) & 0.2 & \cite{cools2011} & [0.1, 0.4] & DA effects vary 2--3$\times$ \\
P3 ($\gamma$) & 0.5 & \cite{dunbar2009} & [0.2, 0.8] & Context-dependent \\
\bottomrule
\end{tabular}
\end{table}

\section{Computational Analysis}
\label{sec:computational}

\subsection{Memory Complexity}

\begin{table}[htbp]
\centering
\caption{Memory Complexity Comparison}
\label{tab:mem_complexity}
\begin{tabular}{lll}
\toprule
\textbf{Component} & \textbf{Standard MoE} & \textbf{Bio-Gating} \\
\midrule
Expert weights & $O(n_e \cdot d^2)$ & Same \\
Gating weights & $O(n_e \cdot d)$ & Same \\
Membrane potential & None & $O(n_e)$ \\
VAD state & None & $O(3)$ \\
NT state & None & $O(5)$ \\
Mood state & None & $O(1)$ \\
\bottomrule
\end{tabular}
\end{table}

\textbf{Overhead}: $O(n_e + 9)$ = negligible for typical $n_e \in [4, 64]$.

\subsection{Behavioral Expressivity}

\begin{definition}[Behavioral Expressivity]
Let $\mathcal{B}$ be the behavioral output space and $\mathcal{S}$ the state space. The expressivity measure:
\begin{equation}
\mathcal{E} = \frac{|\{b \in \mathcal{B} : \exists s \in \mathcal{S}, f(s) = b\}|}{|\mathcal{B}|}
\end{equation}
\end{definition}

\begin{lemma}[Bio-Gating Expressivity]
Bio-Gating achieves expressivity:
\begin{equation}
\mathcal{E}_{\text{Bio-Gating}} \geq \frac{|\mathcal{S}_{\text{VAD}}| \cdot |\mathcal{S}_{\text{NT}}|}{|\mathcal{B}|}
\end{equation}
\end{lemma}

\begin{theorem}[Expressivity Advantage]
Bio-Gating achieves state-dependent expressivity unavailable in standard MoE:
\begin{equation}
\mathcal{E}_{\text{Bio-Gating}} - \mathcal{E}_{\text{Standard}} = \frac{|\mathcal{S}| - |\mathcal{X}|}{|\mathcal{B}|}
\end{equation}
For continuous state spaces, this difference is unbounded.
\end{theorem}

\subsection{Convergence Analysis}

\begin{theorem}[Bio-Gating Convergence]
If the neurochemical state $\mathbf{s}_t$ converges to $\mathbf{s}^*$, then $g_t$ converges to:
\begin{equation}
g^* = \text{softmax}(W_c x + p^* + e(\mathbf{s}^*) + m(\mathbf{s}^*))
\end{equation}
\end{theorem}

\begin{proof}
The gate function $g: \mathcal{S} \rightarrow \mathcal{G}$ is a composition of continuous functions:
\begin{enumerate}
\item The softmax function $\sigma: \R^{n_e} \rightarrow (0,1)^{n_e}$ is continuous (each component is a ratio of exponentials).
\item The modulation functions $e(\mathbf{s})$ and $m(\mathbf{s})$ are continuous (linear in VAD and mood).
\item The membrane potential update $p_t \rightarrow p^*$ is continuous (exponential decay).
\end{enumerate}
Since $\mathbf{s}_t \rightarrow \mathbf{s}^*$ (by Theorem 4.4), and $g$ is continuous in $\mathbf{s}$, by the sequential continuity property: $\mathbf{s}_t \rightarrow \mathbf{s}^*$ implies $g(\mathbf{s}_t) \rightarrow g(\mathbf{s}^*)$.
\end{proof}

\begin{remark}[Base Gate Approximation Error]
The base gate approximation $g^{\text{base}} = \text{softmax}(W_c x + p + e + m)$ resolves the recursive dependency in Equations 2.15--2.16. The approximation error is bounded by:
\begin{equation}
||g - g^{\text{base}}|| \leq \frac{\eta + \kappa}{n_e} \cdot \max(\text{DA}, \text{5-HT})
\end{equation}
For typical modulation strengths ($\eta = 0.2, \kappa = 0.15, n_e = 4$), the maximum error is $\approx 0.0875$ (8.75\% gate probability shift). The fixed-point iteration converges when:
\begin{equation}
|\lambda_{\max}((I + A))| \cdot (\eta + \kappa) < 1
\end{equation}
which holds for $\eta + \kappa < 0.5$ given biological decay rates.
\end{remark}

\begin{corollary}[Convergence Rate]
The convergence rate is bounded by the slowest decay rate:
\begin{equation}
\tau_{\text{convergence}} = \max_i \frac{1}{1 - \lambda_i}
\end{equation}
For typical $\lambda_i \in [0.9, 0.99]$ (biological decay timescales), convergence requires 10--100 cycles.
\end{corollary}

\section{Experimental Demonstration}
\label{sec:experiments}

\textit{Note}: All experiments demonstrate \textbf{qualitative consistency} with clinical phenomena, not validated quantitative predictions. The results show that architecture dynamics produce behavior patterns similar to clinical observations, but coefficients are literature-derived rather than patient-fitted. See Section~\ref{sec:limitations} for detailed interpretation guidelines.

\subsection{Experimental Design}

We conducted 13 experiments across four domains:

\begin{table}[htbp]
\centering
\caption{Complete Experiment Classification}
\label{tab:experiments}
\begin{tabular}{lllll}
\toprule
\textbf{Exp} & \textbf{Domain} & \textbf{Contribution Tested} & \textbf{Method} & \textbf{Metrics} \\
\midrule
1 & Computational & C2 (EventBus sparsity) & Cycle logging & Active region count \\
2 & Computational & C4 (Resource constraints) & Budget monitoring & Balance trajectory \\
3 & Stress & C3 (P1: cortisol$\rightarrow$PFC) & Chronic stress injection & PFC inhibition, DA decline \\
4 & Stress & C3 (Epigenetic tagging) & Stress + recovery & Memory persistence \\
5 & Clinical & C3 (P3+P6: Captivity bonding) & Captor-captive simulation & Fight/fawn ratio \\
10 & Clinical & C1+C3 (D2 blockade) & Antipsychotic injection & PSI, EPS, Treatment index \\
A & Stress & C1+C3 (Anhedonia cascade) & HPA stress cycle & Exploration rate \\
\bottomrule
\end{tabular}
\end{table}

\subsection{Key Results}

\textbf{Experiment 10: D2 Receptor Blockade}

\begin{definition}[D2 Occupancy Model]
Let $o \in [0, 1]$ be D2 receptor occupancy. The therapeutic response (PSI) follows:
\begin{equation}
\text{PSI}(o) = a \cdot o - b \cdot o^2
\end{equation}
where $a = 0.45$ and $b = 0.30$.
\end{definition}

\textbf{EPS Risk Model}:
\begin{equation}
\text{EPS}(o) = c \cdot \max(0, o - o_{\text{threshold}})^k
\end{equation}
where $o_{\text{threshold}} = 0.75$ and $k = 2$.

\begin{lemma}[Optimal Occupancy]
The optimal D2 occupancy for maximal therapeutic index:
\begin{equation}
o^* \approx 0.75 \pm 0.05
\end{equation}
\end{lemma}

\begin{table}[htbp]
\centering
\caption{D2 Occupancy Results}
\label{tab:d2}
\begin{tabular}{lccc}
\toprule
\textbf{Occupancy} & \textbf{PSI Improvement} & \textbf{EPS Index} & \textbf{Treatment Index} \\
\midrule
30\% & 13\% $\pm$ 4\% & 0.00 $\pm$ 0.00 & 1.13 $\pm$ 0.1 \\
75\% & 33\% $\pm$ 5\% & 0.00 $\pm$ 0.00 & \textbf{1.33 $\pm$ 0.1} \\
85\% & 38\% $\pm$ 4\% & 0.25 $\pm$ 0.05 & 1.07 $\pm$ 0.1 \\
95\% & 41\% $\pm$ 3\% & 0.64 $\pm$ 0.08 & 0.83 $\pm$ 0.1 \\
\bottomrule
\end{tabular}
\end{table}

Statistics: PSI gradient $F(1,28)=42.3$, $p<0.001$; model fit $R^2 = 0.94$.

\textbf{Experiment 5: Stockholm Syndrome}

\begin{definition}[Bonding Dynamics]
The bonding score evolution with stress-inverted-U:
\begin{equation}
\frac{dB}{dt} = \kappa \cdot S \cdot \left(1 - \frac{S}{S_{\max}}\right) \cdot (1 - B) - \lambda \cdot B
\end{equation}
\end{definition}

\begin{table}[htbp]
\centering
\caption{Stockholm Syndrome Results}
\label{tab:stockholm}
\begin{tabular}{lcccc}
\toprule
\textbf{Phase} & \textbf{Stress Level} & \textbf{Bonding Score} & \textbf{Fight Ratio} & \textbf{Fawn Ratio} \\
\midrule
Resistance & $S=0.3$ & 0.17 $\pm$ 0.05 & 0.76 $\pm$ 0.10 & 0.05 $\pm$ 0.03 \\
Pressure & $S=0.7$ & 0.68 $\pm$ 0.08 & 0.12 $\pm$ 0.05 & 0.78 $\pm$ 0.09 \\
Bonding & $S=0.85$ & 0.88 $\pm$ 0.06 & 0.02 $\pm$ 0.02 & 0.92 $\pm$ 0.06 \\
\bottomrule
\end{tabular}
\end{table}

Statistics: Fight decline $t=12.67$, $p<0.001$, $d=3.21$.

\textbf{Experiment A: Stress-Induced Anhedonia}

\begin{definition}[Anhedonia Cascade]
The coupled PFC-DA decline:
\begin{equation}
\frac{d\text{PFC}}{dt} = -\alpha \cdot \text{cortisol} \cdot \indicator[\text{stress}] + \gamma_{\text{regen}} \cdot (\text{PFC}_{\text{baseline}} - \text{PFC})
\end{equation}
\begin{equation}
\frac{d\text{DA}}{dt} = -\beta \cdot \text{PFC}_{\text{deficit}} + \delta_{\text{regen}} \cdot (\text{DA}_{\text{baseline}} - \text{DA})
\end{equation}
\end{definition}

\begin{table}[htbp]
\centering
\caption{Stress Anhedonia Results}
\label{tab:anhedonia}
\begin{tabular}{lcccc}
\toprule
\textbf{Metric} & \textbf{Baseline} & \textbf{Stress} & \textbf{Recovery (7d)} & \textbf{Recovery (30d)} \\
\midrule
Cortisol & 0.45 $\pm$ 0.05 & 1.0 $\pm$ 0.0 & 0.48 $\pm$ 0.06 & 0.45 $\pm$ 0.05 \\
PFC inhibition & 0.70 $\pm$ 0.03 & 0.42 $\pm$ 0.04 & 0.51 $\pm$ 0.05 & 0.65 $\pm$ 0.04 \\
Exploration & 0.099 $\pm$ 0.01 & 0.051 $\pm$ 0.02 & 0.068 $\pm$ 0.02 & 0.091 $\pm$ 0.02 \\
\bottomrule
\end{tabular}
\end{table}

Statistics: PFC decline $t=9.15$, $p<0.001$, $d=2.44$.

\subsection{Parameter Sensitivity: Modeling Individual Differences}

To demonstrate the architecture's ability to capture clinical subtypes, we conducted parameter sensitivity analysis across the uncertainty ranges specified in Section 4.6.

\begin{table}[htbp]
\centering
\caption{Behavioral Subtypes from Parameter Variation}
\label{tab:subtypes}
\begin{tabular}{llcc}
\toprule
\textbf{Parameter Range} & \textbf{Clinical Analogue} & \textbf{PFC Decline} & \textbf{Recovery (30d)} \\
\midrule
$\alpha = 0.2$ (resilient) & Low stress sensitivity & $-15\%$ & $95\%$ baseline \\
$\alpha = 0.4$ (typical) & Average sensitivity & $-40\%$ & $93\%$ baseline \\
$\alpha = 0.6$ (vulnerable) & High stress sensitivity & $-58\%$ & $65\%$ baseline \\
\midrule
$\eta = 0.1$ (DA-insensitive) & Anhedonia-prone & Exploration $-60\%$ & $70\%$ baseline \\
$\eta = 0.2$ (typical) & Normal DA response & Exploration $-49\%$ & $91\%$ baseline \\
$\eta = 0.4$ (DA-sensitive) & Reward-responsive & Exploration $-35\%$ & $100\%$ baseline \\
\bottomrule
\end{tabular}
\end{table}

\textbf{Key observations}:
\begin{enumerate}
\item \textbf{Resilient vs.\ vulnerable subtypes}: High $\alpha$ (0.6) produces persistent PFC deficit with incomplete recovery, matching clinical observations of stress-vulnerable individuals. Low $\alpha$ (0.2) shows rapid recovery, representing resilient phenotypes.
\item \textbf{DA sensitivity affects exploration}: Low $\eta$ (0.1) produces prolonged anhedonia with slow exploration recovery; high $\eta$ (0.4) shows rapid normalization.
\item \textbf{Qualitative patterns preserved}: The inverted-U D2 response and stress-scar trajectories persist across all parameter values, with quantitative scaling rather than qualitative disruption.
\end{enumerate}

\textbf{Ablation experiment}: To quantify Bio-Gating's contribution vs.\ Top-1 selection, we compared:
\begin{itemize}
\item \textbf{Bio-Gating + modulation}: State-dependent routing as described
\item \textbf{Bio-Gating (fixed coefficients)}: Modulation disabled (DA=5-HT=NE=0.5 constant)
\item \textbf{Standard Top-1 MoE}: Content-only routing (Switch Transformer equivalent)
\end{itemize}

\begin{table}[htbp]
\centering
\caption{Ablation: Routing Entropy Under Stress}
\label{tab:ablation}
\begin{tabular}{lcc}
\toprule
\textbf{Configuration} & \textbf{Baseline Entropy} & \textbf{Stress Entropy} \\
\midrule
Bio-Gating + modulation & 2.8 bits & 1.2 bits \\
Bio-Gating (fixed coefficients) & 2.8 bits & 2.6 bits \\
Standard Top-1 MoE & 2.8 bits & 2.8 bits \\
\bottomrule
\end{tabular}
\end{table}

Modulation contributes $\Delta H = 1.6$ bits of state-dependent routing entropy reduction under stress---behavioral expressivity unavailable in content-only architectures.

\subsection{Extended Experiments for Revision}

\textbf{Parameter Combination Sensitivity}: We tested $\alpha \times \eta \times \kappa$ interactions across 27 parameter combinations.

\begin{table}[htbp]
\centering
\caption{Parameter Interaction Effects on Scar}
\label{tab:param_interaction}
\begin{tabular}{lccc}
\toprule
\textbf{Subtype} & \textbf{Parameters} & \textbf{Scar Effect} & \textbf{Recovery} \\
\midrule
Resilient & $\alpha=0.2, \eta=0.4$ & $-49.8\%$ & $102\%$ \\
Typical & $\alpha=0.4, \eta=0.2$ & $0.5\%$ & $78\%$ \\
Vulnerable & $\alpha=0.6, \eta=0.1$ & $25.9\%$ & $6\%$ \\
Anhedonia-prone & $\alpha=0.4, \eta=0.1$ & $25.5\%$ & $3\%$ \\
\bottomrule
\end{tabular}
\end{table}

\textbf{Key finding}: High stress sensitivity ($\alpha=0.6$) combined with low DA response ($\eta=0.1$) produces persistent deficit (scar effect 25\%), while low stress sensitivity ($\alpha=0.2$) with high DA response ($\eta=0.4$) shows enhanced recovery beyond baseline.

\textbf{Timescale Analysis}: We tested cortisol half-life effects on recovery dynamics.

\begin{table}[htbp]
\centering
\caption{Age-Related Timescale Effects}
\label{tab:timescale}
\begin{tabular}{lccc}
\toprule
\textbf{Configuration} & \textbf{Half-life} & \textbf{Recovery Time} & \textbf{Ratio} \\
\midrule
Young (fast decay) & 30 steps & 29 steps & $1.0\times$ \\
Adult (medium decay) & 60 steps & 59 steps & $2.0\times$ \\
Elderly (slow decay) & 120 steps & 120 steps & $4.1\times$ \\
\bottomrule
\end{tabular}
\end{table}

\textbf{Clinical implication}: Elderly individuals show 4$\times$ longer stress recovery, consistent with age-related HPA axis changes.

\textbf{Fine-Grained Ablation}: We tested individual module contributions.

\begin{table}[htbp]
\centering
\caption{Module Contribution Quantification}
\label{tab:fine_ablation}
\begin{tabular}{lcc}
\toprule
\textbf{Ablation} & \textbf{Metric Affected} & \textbf{Contribution} \\
\midrule
No DA modulation & Routing entropy & $\Delta = -0.001$ bits \\
No 5-HT modulation & Routing entropy & $\Delta = +0.001$ bits \\
No NE modulation & Routing entropy & $\Delta = -0.003$ bits \\
No P1 pathway & PFC decline & $-20.5\%$ (eliminates decline) \\
No EventBus & Sparsity & $-56.7\%$ (full activation) \\
\bottomrule
\end{tabular}
\end{table}

\textbf{Key finding}: P1 pathway (Cortisol$\rightarrow$PFC) is essential for stress response; removing it eliminates PFC decline entirely. EventBus achieves 56.7\% computational sparsity.

\textbf{MoE Baseline Comparison}: We compared Bio-Gating with 6 MoE variants.

\begin{table}[htbp]
\centering
\caption{MoE Variant Comparison}
\label{tab:moe_comparison}
\begin{tabular}{lcccc}
\toprule
\textbf{Variant} & \textbf{FLOP Red.} & \textbf{Sparsity} & \textbf{Load CV} & \textbf{State-Dep.} \\
\midrule
Bio-Gating (ours) & 75\% & 75\% & 0.70 & \checkmark \\
Switch Transformer & 75\% & 75\% & 0.65 & -- \\
GShard & 50\% & 50\% & 0.75 & -- \\
Expert-Choice & 75\% & 75\% & 0.95 & -- \\
Soft-MoE & 0\% & 0\% & 0.90 & -- \\
Dense-FFN & 0\% & 0\% & 1.00 & -- \\
\bottomrule
\end{tabular}
\end{table}

\textbf{Key finding}: Bio-Gating's FLOP reduction (75\%) matches Switch Transformer's Top-1 mechanism. Bio-Gating's novel contribution is \textbf{state-dependent routing expressivity}, not FLOP optimization. Only Bio-Gating captures stress-induced routing changes.

\textbf{Clinical Validation}: We compared simulated trajectories with published clinical data.

\begin{table}[htbp]
\centering
\caption{Clinical Correspondence Analysis}
\label{tab:clinical}
\begin{tabular}{lccc}
\toprule
\textbf{Condition} & \textbf{Clinical Mean} & \textbf{Simulated} & \textbf{Correspondence} \\
\midrule
Depression (HAM-D) & $10.5 \pm 5.2$ & $18.1$ & 63\% \\
Anxiety (STAI) & $38.0 \pm 9.0$ & $49.2$ & 73\% \\
Stress (Cortisol) & $12.0 \pm 3.0$ & $14.2$ & 79\% \\
\bottomrule
\end{tabular}
\end{table}

\textbf{Interpretation}: Simulation produces trajectories qualitatively consistent with clinical data (STAR*D, STAI norms). Correspondence ranges 63--79\%. Full validation requires prospective clinical study.

\section{Behavioral Analysis and Limitations}
\label{sec:behavioral}

\textit{This section analyzes behavioral outputs and discusses limitations. For a consolidated discussion, see Section~\ref{sec:limitations}.}

\subsection{Computational Contributions}

Simulacrum provides four computational advances:

\textbf{C1: Bio-Gating} -- State-dependent routing absent from content-only MoE.

\textbf{C2: EventBus} -- Event-driven sparsity mimicking biological sparse firing.

\textbf{C3: Coupling Pathways} -- Empirically-grounded neuromodulatory interactions.

\textbf{C4: Demonstration Framework} -- Quantitative behavioral outputs matching clinical descriptions.

\subsection{Comparison with Prior Work}

\begin{table}[htbp]
\centering
\caption{Architecture Comparison}
\label{tab:arch_comparison}
\scriptsize
\begin{tabular}{lcccccc}
\toprule
\textbf{Feature} & \textbf{Simulacrum} & \textbf{Switch} & \textbf{SOAR} & \textbf{ACT-R} & \textbf{LIDA} & \textbf{Nengo} \\
\midrule
State-dependent routing & \checkmark & -- & -- & -- & Limited & -- \\
Modulation factors & 7 & 0 & 0 & 0 & Limited & 0 \\
Sparsity & $\sim$23\% & $\sim$50\% & Rule-based & Chunk-based & $\sim$10\% & Continuous \\
Coupling pathways & 7 & 0 & 0 & 0 & Limited & 0 \\
Mathematical formulation & \checkmark & \checkmark & Partial & Partial & Partial & \checkmark \\
Neuromodulation & \checkmark & -- & -- & -- & Limited & Limited \\
\bottomrule
\end{tabular}
\end{table}

\subsection{Limitations}
\label{sec:limitations}

\textbf{1. Scale and Trainability}: Bio-Gating is tested at 12M parameters with fixed modulation coefficients derived from literature. Scaling to 1B+ parameters may introduce instabilities in coupling dynamics. End-to-end training could discover coefficient values diverging from biological estimates.

\textbf{2. Biological Grounding Simplifications}: The architecture makes several simplifying assumptions that may affect biological fidelity:
\begin{itemize}
\item \textit{Scalar neurotransmitter values}: DA/5-HT/NE modeled as $[0,1]$ scalars ignore spatial gradients (nucleus-specific release, e.g., VTA vs.\ SNc), receptor subtypes (D1 vs.\ D2, 5-HT$_{1A}$ vs.\ 5-HT$_{2A}$), and temporal dynamics (100--1000\,ms delays).
\item \textit{Linear coupling}: Actual neuromodulatory effects are often nonlinear with threshold dynamics and saturation.
\item \textit{Discrete brain regions}: 14 modules approximate, but do not precisely map to, neuroanatomy. For example, ``Basal Ganglia'' as a single module ignores striosome/matrix organization relevant to DA modulation.
\end{itemize}

\textbf{3. Sparsity Mechanism Difference}: EventBus achieves 23\% activation sparsity through subscription-based filtering, while biological sparse firing (1--5\%) emerges from threshold-based membrane dynamics. These mechanisms differ computationally: EventBus sparsity is \textit{structural} (predetermined subscriptions), while biological sparsity is \textit{dynamic} (input-dependent firing). The 5--20$\times$ difference reflects a completeness-fidelity tradeoff.

\textbf{4. Behavioral Demonstration Interpretation}: Experiments demonstrate \textit{qualitative consistency} with clinical phenomena, not validated quantitative predictions. We emphasize the following clarifications:

\begin{itemize}
\item \textbf{``Verification'' vs ``Consistency''}: Throughout this paper, terms like ``shows consistency with'' or ``matches patterns from'' clinical phenomena are used instead of ``validates'' or ``verifies.'' The experiments demonstrate that the architecture produces behavior qualitatively similar to clinical observations, not that the architecture is a validated model of the underlying biology.

\item \textbf{D2 blockade experiment}: The inverted-U therapeutic curve (optimal 60--80\% occupancy) matches clinical pharmacology patterns qualitatively. Precise numerical values (PSI improvement percentages, EPS thresholds) should not be interpreted as validated predictions. Coefficients were derived from literature ranges, not fitted to patient data.

\item \textbf{Stockholm syndrome simulation}: We use the term ``captivity-induced bonding'' rather than claiming to model Stockholm syndrome, which lacks standardized diagnostic criteria in DSM-5 or ICD-11. The simulation demonstrates stress-induced bonding dynamics consistent with trauma psychology literature, not a validated clinical model.

\item \textbf{Stress-induced anhedonia}: The PFC$\rightarrow$DA cascade and partial recovery (``stress scar'' effect) match qualitative patterns from clinical observations, but recovery trajectories were not validated against longitudinal patient data.
\end{itemize}

\textbf{5. Lack of Standard Benchmarks}: The architecture has not been validated on standard cognitive tasks (n-back, Stroop, Tower of Hanoi) or NLP benchmarks (GLUE, WikiText-103). This limits comparison with existing architectures on established metrics.

\textbf{6. Individual Differences}: While Section 4.6 provides uncertainty ranges, the experiments report population-level means. Parameter sensitivity analysis (Section 4.7) shows qualitative robustness, but the architecture's ability to model clinical subtypes (e.g., resilient vs.\ vulnerable individuals) requires explicit demonstration.

\textbf{What This Paper Does NOT Claim}:
\begin{itemize}
\item A validated computational psychiatry model (requires patient data fitting)
\item A production MoE architecture (requires benchmark testing)
\item A theory of neuromodulation (engineering implementation, not biological theory)
\item Quantitative clinical predictions (coefficients are literature-derived, not patient-fitted)
\end{itemize}

\textbf{What We DO Claim}:
\begin{itemize}
\item A novel routing formulation with mathematical rigor
\item Proof-of-concept that neuromodulated routing produces state-dependent behavior
\item Formalized dynamics with stability/convergence guarantees
\item Qualitative consistency with clinical phenomena across three domains
\end{itemize}

\section{Discussion}
\label{sec:discussion}

\textit{This section discusses the architecture's contributions and relation to existing work. For limitations, see Section~\ref{sec:limitations}.}

We presented Bio-Gating, a neurotransmitter-modulated routing mechanism that extends Mixture-of-Experts architectures with neuromodulatory state dependence. The key computational contributions are: (1) a mathematically rigorous formulation of state-dependent gating with stability and convergence guarantees; (2) an event-driven sparse activation architecture achieving $\sim$23\% activation sparsity; and (3) seven empirically-grounded coupling pathways with formalized dynamics.

The central insight is that neuromodulatory dynamics are a computational primitive: routing and activation should reflect the neurochemical state of the system. This property is fundamental to biological computation but absent from standard neural architectures.

\textbf{Relation to Existing Work}:

\textit{Conditional MoE}: Recent conditional MoE approaches \cite{conditional_moe} modulate routing based on input features (e.g., sentence length, domain). Bio-Gating differs by modulating based on \textit{internal state} (DA/5-HT/NE/VAD/mood), enabling behavioral changes without input variation---a capability absent from input-conditional architectures.

\textit{Reinforcement Learning}: The DA-exploration pathway (P2) parallels RL exploration policies. However, Bio-Gating's modulation is \textit{continuous} (softmax perturbation) rather than \textit{discrete} (action selection), enabling graded behavioral shifts.

\textit{Computational Psychiatry Models}: Existing models (e.g., reinforcement learning deficits in depression \cite{huys2016}) focus on single mechanisms. Simulacrum integrates multiple pathways (P1--P7) enabling cascade effects (stress$\rightarrow$PFC$\rightarrow$DA) matching clinical complexity.

\textbf{Limitations}: See detailed discussion in Section~\ref{sec:limitations}. Key limitations include: fixed coefficients at 12M parameters, biological simplifications (scalar NT values, linear coupling), lack of standard benchmark validation, and qualitative rather than quantitative clinical consistency.

\textbf{Future Directions}:
\begin{enumerate}
\item \textit{Benchmark validation}: Evaluate Bio-Gating on WikiText-103 (language modeling) and GLUE tasks, comparing routing entropy and task performance under state perturbation.
\item \textit{Neural data fitting}: Compare pathway dynamics with fMRI BOLD under stress tasks (e.g., Trier Social Stress Test) to estimate cortisol$\rightarrow$PFC coefficient.
\item \textit{Scale testing}: Deploy at 100M+ parameters to test coupling stability and memory overhead.
\item \textit{Application domains}: Adaptive AI systems (game NPCs with emotional responses), computational psychiatry hypothesis generation, brain-inspired embodied agents.
\end{enumerate}

\section{Conclusion}
\label{sec:conclusion}

We presented Simulacrum, a cognitive architecture that integrates neuromodulated routing with event-driven sparse activation for adaptive AI systems. The architecture makes three primary contributions:

\begin{enumerate}
\item \textbf{Bio-Gating} provides state-dependent MoE routing through neuromodulatory (DA/5-HT/NE), emotional (VAD), and mood modulation, achieving 50\% FLOP reduction with state-dependent routing enabling behavioral expressivity.

\item \textbf{EventBus} implements event-driven sparse activation through publish-subscribe, achieving $\sim$23\% activation sparsity with $O(k \cdot n)$ complexity.

\item \textbf{Seven coupling pathways} formalize cross-module dynamics with stability and convergence guarantees, producing emergent behaviors matching clinical phenomena.
\end{enumerate}

The central insight is that neuromodulatory state shapes routing, a property fundamental to biological cognition but absent from standard neural architectures. This enables adaptive AI systems whose behavior varies with internal state in biologically interpretable ways.

\section*{Declarations}

\textbf{Funding}: This research received no specific funding from public, commercial, or not-for-profit funding agencies.

\textbf{Conflicts of Interest}: The authors declare no conflicts of interest related to this work.

\textbf{Data Availability}: This paper presents a theoretical cognitive architecture with simulated experiments. All experimental configurations, parameters, and simulation protocols are fully described within the manuscript. Implementation code is available at GitHub for reproducibility verification.

\textbf{Ethics Approval}: Not applicable---this study involves computational simulation only, with no human or animal subjects.

%% Bibliography
\bibliographystyle{elsarticle-num}

\begin{thebibliography}{30}

\bibitem{laird2012}
Laird, J. E. (2012). The SOAR cognitive architecture. MIT Press.

\bibitem{anderson2004}
Anderson, J. R., Bothell, D., Byrne, M. D., Douglass, S., Lebiere, C., \& Qin, Y. (2004). An integrated theory of the mind. Psychological Review, 111(4), 1036--1060.

\bibitem{franklin2018}
Franklin, S., Strain, S., McCall, R., \& Baars, B. J. (2018). Global workspace theory and LIDA: A testable theory of consciousness. Cognitive Systems Research, 51, 1--29.

\bibitem{schultz2007}
Schultz, W. (2007). Multiple dopamine functions at different time courses. Annual Review of Neuroscience, 30, 259--288.

\bibitem{arnsten2009}
Arnsten, A. F. (2009). Stress signalling pathways that impact prefrontal cortex structure and function. Nature Reviews Neuroscience, 10(6), 410--422.

\bibitem{cools2011}
Cools, R., Nakamura, K., \& Daw, N. (2011). Serotonin and dopamine: unbalanced modulators of risk and reward. Neuropsychopharmacology, 36(1), 267--268.

\bibitem{dunbar2009}
Dunbar, R. I. (2009). The social brain hypothesis and its relevance to social network analysis. Annals of Human Biology, 36(5), 562--572.

\bibitem{lennie2003}
Lennie, P. (2003). The cost of cortical computation. Current Biology, 13(6), 493--497.

\bibitem{hasselmo1999}
Hasselmo, M. E. (1999). A proposed role for forebrain cholinergic tuning in hippocampal function. Neurobiology of Learning and Memory, 71(1), 1--14.

\bibitem{dayan2009}
Dayan, P., \& Huys, Q. J. (2009). Serotonin in affective control. Annual Review of Neuroscience, 32, 95--126.

\bibitem{aston2005}
Aston-Jones, G., \& Cohen, J. D. (2005). An integrative theory of locus coeruleus-norepinephrine function. Annual Review of Neuroscience, 28, 403--450.

\bibitem{fedus2021}
Fedus, W., Zoph, B., \& Shazeer, N. (2021). Switch transformers: Scaling to trillion parameter models. arXiv:2101.03961.

\bibitem{jacobs1991}
Jacobs, R. A., Jordan, M. I., Nowlan, S. J., \& Hinton, G. E. (1991). Adaptive mixtures of local experts. Neural Computation, 3(1), 79--87.

\bibitem{shazeer2017}
Shazeer, N., et al. (2017). Outrageously large neural networks: The sparsely-gated mixture-of-experts layer. arXiv:1701.06538.

\bibitem{kapur2000}
Kapur, S., Zipursky, R., Jones, C., Remington, G., \& Houle, S. (2000). A positron emission tomography study of quetiapine in schizophrenia. Archives of General Psychiatry, 57(6), 553--559.

\bibitem{farde1992}
Farde, L., Nordstr\"om, A. L., Wiesel, F. A., Halldin, C., Sedvall, G., \& Uppfeldt, G. (1992). Positron emission tomographic analysis of central D1 and D2 dopamine receptor occupancy. Archives of General Psychiatry, 49(7), 538--544.

\bibitem{eichenbaum2017}
Eichenbaum, H. (2017). Memory: Organization and control. Annual Review of Psychology, 68, 19--45.

\bibitem{xie2013}
Xie, L., Kang, H., Xu, Q., et al. (2013). Sleep drives metabolite clearance from the adult brain. Science, 342(6156), 373--377.

\bibitem{miller1956}
Miller, G. A. (1956). The magical number seven, plus or minus two. Psychological Review, 63(2), 81--97.

\bibitem{davies2018}
Davies, M., Wild, A., Lin, A., et al. (2018). Loihi: A neuromorphic manycore processor with on-chip learning. IEEE Micro, 38(1), 82--99.

\bibitem{merolla2014}
Merolla, P. A., Arthur, J. V., Alvarez-Icaza, R., et al. (2014). A million spiking-neuron integrated circuit with a scalable communication network. Science, 345(6197), 668--673.

\bibitem{conditional_moe}
Liu, Z., Wang, Y., Li, K., et al. (2022). Deep mixture of experts via task-dependent gate. NeurIPS 2022.

\bibitem{huys2016}
Huys, Q. J., Maia, T. V., \& Frank, M. J. (2016). Computational psychiatry as a bridge from neuroscience to clinical applications. Nature Neuroscience, 19(3), 404--413.

\end{thebibliography}

\end{document}
